\chapter{Kết luận}
\label{ch:ketluan}

\section{Đóng góp của luận văn}

Đề tài xuất phát từ một quan sát đơn giản. Cho cùng một ảnh màn hình và cùng một mục
tiêu, dòng nghiên cứu tác tử giao diện sinh ra toạ độ để máy tự bấm; người dùng đang bí
giữa chừng lại cần một câu tiếng Anh nói cho họ biết phải chạm vào đâu. Khi đổi đầu ra
như vậy, khâu sinh câu không phải phần khó nhất. Phần khó nằm ở chỗ không có cách nào
chấm được câu sinh ra, vì cùng một nút có nhiều tên gọi đều đúng và không bộ dữ liệu nào
liệt kê hết. Vì lý do đó, phần lớn công sức của luận văn dồn vào việc dựng và kiểm chứng
một dụng cụ đo.

\subsection*{Về đo lường}

Thước \textbf{executability} đưa câu sinh ra cho một mô hình định vị phần tử độc lập,
không hề thấy toạ độ đáp án, rồi hỏi xem điểm nó trả về có rơi vào ô Voronoi của phần tử
người dùng đã chạm hay không. Điều đáng nói ở đây không phải bản thân ý tưởng, vốn có gốc
từ dòng sinh biểu thức quy chiếu, mà là khối lượng đo đạc đặt bên dưới nó:

\begin{itemize}[leftmargin=1.4em,itemsep=0.2em]
\item Một cổng chặn về sai số của dụng cụ đo, khoá lại trước khi tiêu tài nguyên và suy
ra từ hình học màn hình. Kết quả đo cho sai số trung vị $0{,}7\%$ bề ngang màn so với
ngưỡng $3\%$.
\item Luật xác định trúng được chọn từ \emph{sàn} mà nó đạt được. Đặt trước một câu gọi
tên sai nút, ngưỡng dung sai quy ước vẫn cho điểm $84{,}3\%$ số ca, để lại biên độ động
quá hẹp; luật ô Voronoi hạ con số đó xuống $2{,}8\%$.
\item Trần của thước là một đại lượng \emph{đo được}: chấm chính câu do người viết trên
đủ $4.463$ bước cho $75{,}7\%$, KTC $[74{,}1; 77{,}3]$. Mọi điểm số khác đọc trên nền này.
\item Mười phép bơm lỗi với ngưỡng đóng băng trước khi chạy, và một phân tích chỉ ra
phần lớn trong đó không thể trượt theo cấu tạo (Mục~\ref{sec:bomloi}).
\item Một phép thử độ bền: chấm lại ba nhánh dưới năm luật khác nhau. Mức tuyệt đối di
chuyển tới $26$ điểm, song thứ tự giữa các hệ thống giữ nguyên ở cả năm luật.
\end{itemize}

\subsection*{Về dữ liệu và mô hình}

Tập giám sát $64.567$ bước được dựng hoàn toàn bằng máy, không thuê người dán nhãn và
không đưa nhãn do một mô hình ngôn ngữ khác sinh ra vào vòng huấn luyện. Ba phép kiểm bảo
vệ phần này: phép ghép đa nguồn được đối chiếu với một đối chứng lệch chủ ý ($48\%$ so
với $20\%$), chín bất biến cấu trúc đều đạt ở quy mô đủ, và không có tác vụ nào trùng
giữa tập dạy với tập kiểm. Thành phần đề xuất chạy trên nền dữ liệu đó chưa cho kết quả tại thời điểm
nộp; thứ đã kiểm được là chính đường ống, không phải hiệu lực của thành phần.

\subsection*{Về kết quả định lượng}

Trên $4.463$ bước chạm, mô hình Qwen2.5-VL-3B sau tinh chỉnh QLoRA đạt $59{,}1\%$, trong
khi chính mô hình đó khi chưa tinh chỉnh chỉ đạt $47{,}6\%$; cả hai đọc trên nền trần
$75{,}7\%$. Chênh lệch ghép cặp $11{,}5$ điểm ($\chi^2 = 243{,}2$, $p < 0{,}001$) đứng
vững qua \emph{năm} cách giải thích thay thế; cách thứ sáu, bắt chước văn phong, chưa
loại trừ được và nó chồng lấn với đường nhiễm của dụng cụ đo nêu ở Mục~\ref{sec:hanche}.
Trong năm phép đã bác được có một phép mà thước dựa trên câu chuẩn không tự làm giúp: loại bỏ toàn bộ những câu mô hình chép lại nguyên văn
câu chuẩn, và sau khi loại thì khoảng cách vẫn còn $8{,}4$ điểm. Ba khoảng tin cậy rời
nhau, cho thấy dụng cụ đo phân giải được hai hệ thống thật.

\subsection*{Về quy trình}

Thiết kế so sánh, cùng luật đọc cho cả bốn kết cục có thể xảy ra, được đưa vào kho phiên
bản trước lượt huấn luyện đầu tiên. Các mục sửa đổi được ghi kèm ngày tháng ở cuối bản đăng ký, phần
lớn ra đời khi chưa tồn tại bất kỳ điểm số nào. Những con số bị rút trong quá
trình làm cũng được giữ lại cùng lý do rút, chẳng hạn ước lượng trần $70{,}0\%$ tính trên
mẫu con $300$ bước đã bị thay bởi $75{,}7\%$ đo trên toàn tập.

\subsection*{Về mức độ hoàn thành}

Đại lượng chính của thiết kế, $\Delta = \text{S2} - \text{S1}$, chưa được báo cáo vì
nhánh S2 chưa chạy xong tại thời điểm nộp bản này (Mục~\ref{sec:chuachay}). Phép so S1
với Base, tuy đã có số, không được dùng để thế chỗ cho nó.

\section{Hạn chế}
\label{sec:hanche}

Phần này liệt kê những chỗ mà kết luận của luận văn còn hở, kèm số đo cho từng chỗ khi có
thể đo được.

\subsection{Hạn chế của dụng cụ đo}

\begin{itemize}[leftmargin=1.4em,itemsep=0.35em]

\item \textbf{Chỉ một mô hình định vị, và còn một đường nhiễm chưa đóng được.} Mọi con số
đều có điều kiện trên UGround-V1-2B. Mô hình này dựng trên Qwen2-VL, cùng dòng với
Qwen2.5-VL-3B đang bị chấm, và bảng thành phần dữ liệu của nó liệt kê $47$ nghìn phần tử
có nhãn người lấy từ split train của AndroidControl. Ở mức màn hình thì tập kiểm vẫn tách
rời, nên đây không phải rò rỉ nhãn. Vấn đề nằm chỗ khác: dụng cụ đo đã quen \emph{văn
phong chú thích} của bộ ngữ liệu này, trong khi S1 lại được dạy viết đúng văn phong ấy,
và một dụng cụ giải mã văn phong đó dễ hơn sẽ đẩy điểm S1 lên. \textbf{Không phép kiểm
nào trong sáu phép ở Mục~\ref{sec:phanbien} loại trừ được cách đọc này.} Chỉ có một phép
đo đóng được nó, là chấm lại bằng một dụng cụ đã xác minh sạch AndroidControl, và phép đó
đã đăng ký làm bước kế tiếp. Có một lập luận giảm nhẹ, tuy chưa đủ để bác bỏ: nếu cơ chế
nhiễm là thật thì nó tác động mạnh nhất lên nhánh câu người viết, nhánh bám văn phong ấy
sát nhất, nên nó sẽ nâng trần trước khi nâng bất kỳ chênh lệch nào.

\item \textbf{Chưa trả lời được câu hỏi về độ mong manh trước cách diễn đạt.} Jandial và
cộng sự~\cite{groundersunderstand} báo cáo các mô hình định vị giao diện trượt tới $84\%$
số ca khi chỉ đổi cách nói, kết quả nhắm thẳng vào thước ở đây. Bộ bơm lỗi
(Mục~\ref{sec:bomloi}) không phản bác được, vì nó nạp một điểm tổng hợp thẳng vào hàm
trúng và \emph{không gọi mô hình định vị lần nào}. Dòng đầu của Bảng~\ref{tab:bomloi} chỉ
cho biết các cổng về chữ không bác nhầm một câu diễn đạt khác; nó im lặng về việc mô hình
định vị có còn tới đúng phần tử hay không khi cách nói đổi đi.

\item \textbf{Vùng mù theo cả hai chiều.} Với độ lệch dưới khoảng $63$ điểm ảnh thì thước
không phát hiện được. Nói chính xác, nó đo được việc gọi tên sai phần tử và bỏ qua việc
trỏ hơi lệch bên trong đúng phần tử. Theo chiều ngược lại, $1.083$ bước mà câu người viết
cũng trượt phần lớn là giới hạn của dụng cụ, và có $935$ bước ($21\%$ quần thể) đánh bại
cả ba nhánh. Trần $75{,}7\%$ vì vậy phản ánh năng lực của dụng cụ đo hơn là năng lực của
ngôn ngữ.

\item \textbf{Luật ô Voronoi chỉ chặt bằng độ đầy đủ của bộ liệt kê phần tử.} Cây trợ
năng bỏ sót một nút biểu tượng cạnh phần tử đích thì một câu sai vẫn lọt. Chiều ngược
lại, các hộp trùng lặp, đã được xử bằng phép gộp theo bán kính, nhưng phép gộp ấy xoá
nhầm một phần tử thật sự riêng biệt ở $8{,}2\%$ số bước.

\item \textbf{Khái niệm được đo hẹp hơn khái niệm người đọc quan tâm.} Thước hỏi một
\emph{mô hình} định vị có tới đúng phần tử từ câu hay không. Câu hỏi mà người dùng quan
tâm là một \emph{người} có tới được hay không. Dòng sinh biểu thức quy chiếu vẫn coi
thành công của người nghe là bằng chứng về chất lượng mô
tả~\cite{mao2016,luo2017,yu2017}, và barem ở đây là toạ độ một người thật đã chạm, song
hai đại lượng ấy không đồng nhất và mức đồng thuận giữa chúng \emph{chưa được đo}. Phép
chấm tay $100$ câu bởi hai người đã đăng ký nhưng chưa chạy.

\end{itemize}

\subsection{Hạn chế của dữ liệu và thiết kế}

\begin{itemize}[leftmargin=1.4em,itemsep=0.35em]

\item \textbf{Tập kiểm không phải ``ứng dụng chưa từng thấy''.} $95{,}6\%$ số bước quy được tên
thuộc về ứng dụng cũng xuất hiện ở tập dạy ($91{,}8\%$ nếu đếm theo ứng dụng riêng biệt), nên mọi phép so với một mô hình bên ngoài chỉ có giá
trị tham khảo. Đo trên chính mô hình của đề tài thì ảnh hưởng không xuất hiện
(Bảng~\ref{tab:sannha}); với $78$ bước ở nhóm chưa thấy, kết luận đúng phải là
\emph{không phát hiện được ảnh hưởng}, một phát biểu yếu hơn hẳn \emph{đã chặn được ảnh
hưởng}.

\item \textbf{Hơn một nửa số bước không quy được về ứng dụng.} $3.828$ trong $6.958$ bước
mang nhãn \emph{không biết}. Nhãn này không bao giờ được đọc thành \emph{chưa thấy}.

\item \textbf{Tách theo tác vụ không đồng nghĩa tách theo cách diễn đạt.} Bộ ngữ liệu
dùng đi dùng lại các cách nói phổ biến, đến mức $17{,}3\%$ câu chuẩn của tập kiểm đã xuất
hiện nguyên văn ngay trong một lát chỉ chiếm $2{,}6\%$ tập dạy.

\item \textbf{Nhãn mô tả còn nhiễu.} $24{,}1\%$ ô dấu hiệu phân biệt không phân giải được
gì, và $22{,}0\%$ phần tử đích không có tên nào để điền vào ô tên.

\item \textbf{Khâu chấm là teacher-forced trên ngữ cảnh.} Trường lịch sử trong câu nhắc
chứa câu người viết ở các bước trước, không phải nhật ký thao tác của mô hình. Vì là hằng
số của thí nghiệm nên nó không giải thích được chênh lệch giữa các nhánh, nhưng mọi con
số tuyệt đối phải đọc kèm điều kiện này.

\item \textbf{Một hạt giống, một lượt huấn luyện.} Con số $11{,}5$ điểm gắn với
\emph{lượt huấn luyện này}. Chưa có hạt giống thứ hai nên chưa biết hai lượt chạy của
cùng một thiết kế lệch nhau bao nhiêu.

\item \textbf{Tầng mức hai đã đăng ký nhưng chưa cài đặt.} Hàm mất mát có khoản phạt cho
câu mơ hồ chưa được viết. Trạng thái đúng của nó là \emph{đã đăng ký, không thực hiện}.

\item \textbf{Phạm vi chung.} Một màn hình, giao diện tiếng Anh, một bộ dữ liệu, một mô
hình định vị, một mô hình gốc cỡ ba tỉ tham số.

\end{itemize}

\section{Hướng phát triển}

Các hướng dưới đây xếp theo mức độ cấp thiết đối với độ tin cậy của kết luận, không xếp
theo độ hấp dẫn về mặt kỹ thuật.

\begin{itemize}[leftmargin=1.4em,itemsep=0.35em]

\item \textbf{Hoàn tất thiết kế đã đăng ký.} Chạy nốt hạt giống thứ hai của S1 để có sàn
nhiễu, khoá ngưỡng, rồi mới huấn luyện S2 cùng hai nhánh đối chứng S2r và S2-nopoint.
Trình tự này đã khoá và không được đảo.

\item \textbf{Kiểm chéo bằng một dụng cụ đo thứ hai.} Về độ tin cậy thì đây là việc quan
trọng nhất, bởi nó là phép đo duy nhất đóng được đường nhiễm nêu ở trên. Ứng viên đã tra
và chọn là một mô hình định vị mã nguồn mở, đạt kết quả cao hơn trên bộ chuẩn di động và
đã xác minh không dùng AndroidControl khi huấn luyện. Chi phí ước tính hơn một giờ máy
cho một lát khoảng $500$ bước trên hai nhánh.

\item \textbf{Đo độ mong manh trước cách diễn đạt.} Viết lại câu chuẩn theo bốn kiểu
(đổi động từ, đảo trật tự, bỏ từ chỉ vị trí, kết hợp cả ba) rồi chạy mô hình định vị trên
các câu đã viết lại. Chỉ tốn một lượt chấm, và trả lời trực tiếp cảnh báo
của~\cite{groundersunderstand}.

\item \textbf{Nối thước với đánh giá của người.} Chấm tay $100$ câu bởi hai người độc lập
và báo cáo mức đồng thuận. Đây là điều kiện để phát biểu mạnh hơn về giá trị cấu trúc của
thước.

\item \textbf{Cài đặt tầng mức hai.} Viết hàm mất mát có khoản phạt lề dựa trên trường mô
tả hàng xóm đã dựng sẵn, thử trên một lát nhỏ trước khi mở rộng. Đây là chỗ lấp trực tiếp
khe hở của Widget Captioning~\cite{widgetcap}.

\item \textbf{Cải thiện bộ liệt kê phần tử.} Vì độ chặt của luật ô Voronoi bị chặn bởi độ
đầy đủ của bộ liệt kê, việc ghép thêm một bộ phát hiện phần tử dựa trên ảnh bên cạnh cây
trợ năng sẽ nâng trực tiếp độ chặt của thước, rõ nhất ở các nút chỉ có biểu tượng.

\item \textbf{Mở rộng phạm vi.} Ba hướng theo chi phí tăng dần: chuỗi nhiều màn hình liên
tiếp; giao diện tiếng Việt, ít nhất ở mức đánh giá định tính; và một bộ dữ liệu thứ hai
để tách kết luận khỏi đặc thù của AndroidControl.

\end{itemize}

\section{Tính sẵn có của dữ liệu và mã nguồn}

Toàn bộ số liệu trong luận văn tái tạo được từ mã nguồn và tệp kết quả thô đi kèm.

\textbf{Dữ liệu nguồn} là AndroidControl~\cite{androidcontrol}, giấy phép CC0, lấy qua
hai bản phân phối công khai trên Hugging Face đã nêu ở Mục 3.1. Ảnh màn hình không được
phân phối lại.

\textbf{Mã nguồn} gồm các bước dựng dữ liệu (ghép nguồn, OCR, dựng nhãn mô tả, dựng bốn
nhánh), tệp cấu hình huấn luyện dùng chung cho mọi nhánh, chương trình sinh câu, và
chương trình chấm điểm cùng phần tự kiểm chạy được của luật chấm.

\textbf{Kết quả thô} được lưu ở mức từng bước. Mỗi nhánh có một tệp câu sinh ra, một tệp
điểm tổng hợp, và một tệp ghi mỗi bước một dòng kèm toạ độ mô hình định vị trả về, toạ độ
đáp án, số phần tử trên màn và giá trị của cả ba điều kiện chấm. Nhờ tệp cuối, việc đổi
luật chấm hay thêm một lát cắt phân tích chỉ cần tính lại từ đó, khỏi phải gọi lại mô
hình định vị, tiết kiệm khoảng $5{,}6$ giờ máy mỗi lần.

Bản đăng ký trước cùng toàn bộ mục sửa đổi có ngày nằm trong kho phiên bản của đề tài,
kèm danh sách các con số đã bị rút với lý do rút và giá trị thay thế.
